\section{Conclusion}
\label{chap:conclusion}

This paper presented a rigorous investigation into the performance and energy efficiency of three BP-free training algorithms: FF, CaFo, and MF, benchmarked against standard backpropagation. Through a framework of fair comparisons on native architectures, incorporating direct hardware measurements, systematic hyperparameter optimization, and validation-based early stopping, this work clarifies the practical viability of these alternatives for more sustainable deep learning. The evolutionary trajectory from FF's foundational proof of concept, through CaFo's structured design, to MF's emergence as a highly efficient and performant method for MLPs highlights the significant potential and rapid advancement within this domain.

\subsection{Key Findings and Answers to Research Questions}
\label{sec:conclusion_findings}

The experimental results provide definitive answers to the research questions posed at the beginning of this work. Regarding classification performance, the MF algorithm, in its native MLP architectures, \textbf{robustly and consistently matched or surpassed the accuracy of its BP baseline}. Its convergence to a lower final validation loss indicates that its layer-wise optimization strategy identified a more effective solution. In contrast, FF achieved comparable accuracy but with markedly slower and more volatile convergence, while the CaFo-Rand-CE variant revealed a clear performance deficit on complex datasets. The CaFo-DFA-CE variant dramatically narrowed this gap, but at a significant computational cost during pre-training.

In terms of direct energy consumption, the MF algorithm yielded significant savings, consuming approximately 41\% less energy than BP on the CIFAR-10 MLP task while achieving superior accuracy. This is in stark contrast to the case for FF, which consumed substantially more energy because of its inefficiency, and CaFo-DFA-CE, which was also more energy-intensive. This research also \textbf{empirically refutes the generalized assumption that all BP-free methods are inherently more memory efficient}. No algorithm yielded dramatic memory savings; for MF on larger MLPs, the peak memory advantage was modest (4\% to 5\%), as the overhead for its additional projection matrices was found to partially offset the theoretical gains.

These results clarify the different energy-performance trade-offs. FF offers an unfavorable trade-off of competitive accuracy for prohibitive inefficiency. CaFo presents a stark choice between the Rand-CE variant's modest efficiency gains for a significant accuracy drop, and the DFA-CE variant's competitive accuracy for a substantial efficiency cost. Crucially, MF delivered the most advantageous trade-off, achieving \textbf{accuracy parity or superiority with simultaneous and significant improvements in training time and energy efficiency}. It also demonstrated effective scalability in both performance and efficiency across the tested MLP and dataset complexities. In summary, this research established FF as a foundational proof of concept, identified CaFo as a structured method defined by clear trade-offs, and distinguished MF as a highly promising algorithm that, for MLP architectures, fulfills the objectives of achieving both BP-competitive accuracy and superior training efficiency.

\subsection{Summary of Contributions}
\label{sec:conclusion_contributions}

This paper makes several key contributions to the study of energy-efficient deep learning. First, it established a \textbf{rigorous and fair benchmarking framework} by replicating the native architectures for FF, CaFo, and MF, creating identical BP baselines, and applying systematic, validation-based optimization to all methods. This work presented one of the first \textbf{extensive independent evaluations of the recently proposed MF algorithm} on its native MLP architectures, robustly demonstrating its ability to improve both accuracy and efficiency. Through detailed profiling centered on \textbf{direct measurements of GPU energy consumption via NVML}, this research provided critical insights into the inefficient hardware utilization of FF and the nuanced memory characteristics of MF.
\pagebreak

This research systematically detailed the trade-offs for each algorithm, clarifying that MF (on MLPs) provided the most favorable balance. Finally, to promote reproducibility and facilitate future research, all developed code, structured configurations, experimental protocols, and measurement tools were organized into an open research infrastructure and made publicly available.

\subsection{Future Research Directions}
\label{sec:conclusion_future_work}

The findings of this paper suggest several promising avenues for future research. A primary direction is architectural exploration, particularly investigating \textbf{MF's adaptability and performance on more complex architectures} like CNNs and Transformers, which will require adapting its projection matrix mechanism. For CaFo, developing more computationally efficient, BP-free methods for training its neural blocks is essential. Further research should also evaluate the most promising algorithms on larger-scale datasets like ImageNet to fully ascertain their scalability. Other important frontiers include developing custom hyperparameter pruners for layer-wise training, deepening the theoretical understanding of how MF's local optimization can achieve superior global solutions, and pursuing \textbf{hardware co-design for BP-free algorithms} to unlock additional energy savings. Extending these methods beyond classification to domains like generative modeling and reinforcement learning, as well as investigating their robustness and uncertainty calibration, will be crucial for their broader adoption.

\subsection{Broader Implications}
\label{sec:conclusion_implications}

The research detailed in this paper has broad implications for the field. As artificial intelligence becomes increasingly prevalent, the energy consumption of these systems has become a significant social concern. The development and adoption of algorithms such as MF, which can significantly reduce the footprint of training energy without a performance penalty, represent \textbf{critical steps toward a more sustainable future for AI}. Lowering energy requirements can also \textbf{democratize AI development}, making it more accessible to researchers with limited computational resources, and is paramount for enabling powerful on-device learning under strict power constraints. The biological inspiration underlying these algorithms fosters a symbiotic relationship between AI and neuroscience, where brain-inspired computation can inform more efficient AI. Ultimately, this research contributes to a paradigm shift in which efficiency is treated as a primary design criterion, on par with accuracy. The journey beyond backpropagation is not just about devising novel algorithms; it is about \textbf{shaping a more efficient, accessible, and responsible future for artificial intelligence}.